\section{Wire and transcript profile v1}
\label{sec:wire-profile}

\begingroup\small

\begin{normative}
This section fixes the bytes consumed by a conforming prover and verifier.
The profile has two physical proof streams, prefix-free typed records, a
BLAKE3 transcript, SHA-256 Merkle roots, and the recursive Ligerito and
proof-of-work parameters below.  A proof using the legacy positional/bincode
codec is not a v1 proof.  End-to-end use additionally requires the three
statement-bound subordinate identifiers to resolve as specified below.
\end{normative}

\subsection{Profile identity and fixed algorithms}
\label{sec:wire-profile-id}

The eight-byte container magic is the ASCII string \code{F2ZPCS\textbackslash0\textbackslash0},
the wire version is one, and the 32-byte profile identifier is
\begin{equation}
 \ProfileID=\operatorname{ASCII}
   (\code{F2ZPCS-v1-B3-SHA256-CORE-0000001}).
 \label{eq:wire-profile-id}
\end{equation}
The literal in Equation~\eqref{eq:wire-profile-id} has exactly 32 bytes and
hex encoding
\code{46325a5043532d76312d42332d5348413235362d434f52452d30303030303031}.
Changing any normative rule in this section requires a new profile identifier.

The transcript and its XOF are unkeyed BLAKE3; Merkle nodes and roots use
SHA-256; PoW uses unkeyed BLAKE3-256.  The field basis is the polynomial basis
of Equation~\eqref{eq:concrete-domains}.  An application PIOP, a GKR
implementation, and the constrained-code construction each supply a distinct
32-byte subordinate profile identifier.  The public statement binds all three
identifiers.  An unknown identifier, an application/GKR profile without a
byte-exact event ledger, or a code profile without a canonical codeword
construction and root vectors is unsupported and \must{} be rejected.

For the deployed range $22\le m\le35$, let $J_m$ be the number of recursive
opening levels.  The level shapes and all parameters except the base fold
difficulties are fixed by
\begin{equation}
 \begin{aligned}
  h_i&:=m-11-3i,&
  J_m&:=\min\{J\ge1:h_{J-1}\in\{3,4,5\}\},\\
  r_i&:=1+i,&
  k_i&:=\begin{cases}4,&i=0,\\3,&i>0,\end{cases}\\
  Q_i&:=(183,90,60,45,36,31,27,24)_i,&
  g_i&:=16,\\
  o_i&:=\begin{cases}0,&i=0,\\1,&i>0.\end{cases}
 \end{aligned}
 \label{eq:ligerito-v1-shape}
\end{equation}
Here $h_i$ is the message-column exponent, $r_i$ the inverse-rate exponent,
$k_i$ the fold count, $Q_i$ the query count, $g_i$ the query-grinding
difficulty, and $o_i$ the OOD sample count.  Implementations \must{} use the
integers in Equation~\eqref{eq:ligerito-v1-shape}, not a floating-point
recomputation.

The remaining base fold difficulties $f_i$ are the following literal table.
Only the first $J_m$ entries of the row for $m$ exist.

\begingroup\footnotesize
\renewcommand{\arraystretch}{1.02}
\begin{longtable}{@{}c@{\quad}l@{}}
\toprule
$m$&$(f_0,\ldots,f_{J_m-1})$\\
\midrule
22&$(9,6,3)$\\
23&$(10,7,4,1)$\\
24&$(11,8,5,2)$\\
25&$(12,9,6,3)$\\
26&$(13,10,7,4,2)$\\
27&$(14,11,8,5,3)$\\
28&$(15,12,9,6,4)$\\
29&$(16,13,10,7,5,3)$\\
30&$(17,14,11,8,6,4)$\\
31&$(18,15,12,9,7,5)$\\
32&$(19,16,13,10,8,6,3)$\\
33&$(20,17,14,11,9,7,4)$\\
34&$(21,18,15,12,10,8,5)$\\
35&$(22,19,16,13,11,9,6,6)$\\
\bottomrule
\end{longtable}
\endgroup

\subsection{Primitive canonical encodings}
\label{sec:wire-primitives}

All unsigned integers are fixed-width little-endian.  Decoders \mustnot{}
reduce, truncate, ignore high bits, or convert to a host-sized integer before
checking that the value fits.  The following encodings are the only canonical
ones.

\begingroup\footnotesize
\renewcommand{\arraystretch}{1.05}
\begin{tabularx}{\linewidth}{@{}p{1.25in}p{1.35in}X@{}}
\toprule
\textbf{Type}&\textbf{Bytes}&\textbf{Canonicality rule}\\
\midrule
$u_8,u_{16},u_{32},u_{64},u_{128}$&1, 2, 4, 8, 16&Unsigned
little-endian of exactly the named width.\\
$a\in\K$&16&$\LE{64}{a_{\lo}}\parallel\LE{64}{a_{\hi}}$ in the fixed
polynomial basis; coefficient bit $j$ is bit $j$ of the resulting 128-bit
word.\\
$a\in\Fq$&16&$\LE{128}{\operatorname{can}_q(a)}$; reject an integer at least
$q=2^{100}-15$.\\
Chunk exponent $u$&16&$\LE{128}{u}$; require $u<2^{127}$.\\
Nonce&8&$\LE{64}{n}$.\\
Merkle value&32&Raw SHA-256 digest bytes, in digest output order.\\
Vector $[x_0,\ldots,x_{n-1}]$&4 plus elements&$\LE{32}{n}$ followed by
canonical elements in index order; $n$ must equal the profile-derived expected
count.\\
Opened rows and path&nested&Vector of rows, each a vector of $\K$ elements,
followed by a vector of 32-byte Merkle values.  Expected row, width, and path
counts are derived before allocation.\\
\bottomrule
\end{tabularx}
\endgroup

Every composite is concatenation without padding.  Options, Rust enums,
\code{usize}, varints, serde defaults, and bincode layouts are not wire types.
For every accepted primitive or composite $x$, decoding its encoding must
consume the whole bounded slice and re-encoding the result must reproduce the
same bytes.

\subsection{Merkle rows and canonical multiproofs}
\label{sec:wire-merkle}

For opening level $0\le k<J_m$, define the message dimension, tree depth,
leaf count, and row width by
\begin{equation}
 n_k:=h_k+k_k,\qquad D_k:=h_k+r_k,\qquad
 \Lambda_k:=2^{D_k},\qquad I_k:=2^{k_k}.
 \label{eq:merkle-level-geometry}
\end{equation}
Thus $n_0=m-7$, $n_k=h_{k-1}$ for $k>0$, $D_k=m-10-2k$,
$I_0=16$, and $I_k=8$ for $k>0$.  Let
$C_k\in\K^{\Lambda_k\times I_k}$ be the level-$k$ codeword matrix, with
leaves and lanes in increasing integer order.  Its exact row and tree bytes
are
\begin{equation}
 \begin{aligned}
  \mathsf{RowBytes}_k(q)
   &:=\operatorname{Enc}_{\K}(C_k[q,0])\parallel\cdots\parallel
      \operatorname{Enc}_{\K}(C_k[q,I_k-1]),\\
  H_{k,0,q}&:=\operatorname{SHA256}(\mathsf{RowBytes}_k(q)),\\
  H_{k,d+1,j}&:=\operatorname{SHA256}
      (H_{k,d,2j}\parallel H_{k,d,2j+1}),\\
  \mathsf{root}_k&:=H_{k,D_k,0}.
 \end{aligned}
 \label{eq:merkle-v1-tree}
\end{equation}
Here $\operatorname{Enc}_{\K}$ is the 16-byte encoding in
Section~\ref{sec:wire-primitives}.  Equation~\eqref{eq:merkle-v1-tree} uses
no native-memory cast, row-length prefix, row index, level tag, padding, or
leaf/internal-node domain byte.  Adding such a byte changes every root and
requires a new profile identifier.  The constrained-code profile defines
how $C_k$ is generated; this section fixes its portable row and tree encoding.

Let the already-derived query set be
$A_0=(q_0<\cdots<q_{Q_k-1})$.  At depth $d$, scan the active positions $A_d$
in increasing order.  If both siblings are active, combine the pair once and
emit no digest.  Otherwise emit the digest at position
$p\mathbin{\mathsf{xor}}1$, using $p$'s parity to place it on the left or
right.  Then set
$A_{d+1}:=\operatorname{sortuniq}\{\lfloor p/2\rfloor:p\in A_d\}$.
The canonical multiproof is the concatenation of emitted digests, bottom-up
and in that scan order, and its exact count is
\begin{equation}
 M_k:=\sum_{d=0}^{D_k-1}
 \left|\{p\in A_d:(p\mathbin{\mathsf{xor}}1)\notin A_d\}\right|.
 \label{eq:merkle-proof-count}
\end{equation}
The level-$k$ opening payload is exactly
\begin{equation}
 \LE{32}{Q_k}\parallel
 \prod_{j=0}^{Q_k-1}
   \bigl(\LE{32}{I_k}\parallel\mathsf{RowBytes}_k(q_j)\bigr)
 \parallel\LE{32}{M_k}\parallel
 \prod_{u=0}^{M_k-1}\mathsf{path}_k[u],
 \label{eq:merkle-opening-payload}
\end{equation}
with byte length
\begin{equation}
 8+Q_k(4+16I_k)+32M_k.
 \label{eq:merkle-opening-length}
\end{equation}
The verifier derives $A_d$, $M_k$, and
Equation~\eqref{eq:merkle-opening-length} from the sampled queries before
allocating or reading the bounded Hint record, requires rows in $q_j$ order,
reconstructs Equation~\eqref{eq:merkle-v1-tree}, and rejects unused path
digests or payload bytes.

\subsection{Outer two-stream container}
\label{sec:wire-container}

The 72-byte header is fixed as follows.  Offsets and sizes are decimal; all
multibyte integers use Section~\ref{sec:wire-primitives}.

\begin{lstlisting}[style=wire]
offset  size  field
0       8     magic = 46 32 5a 50 43 53 00 00
8       2     wire_version = 01 00
10      2     header_len = 48 00                 (72)
12      4     flags = 00 00 00 00
16      32    profile_id
48      4     narg_record_count
52      4     hint_record_count
56      8     narg_byte_len
64      8     hint_byte_len
72      ...   NARG records, then Hint records
\end{lstlisting}

Writing $\mathsf{Hdr}_{72}$ for those exact header bytes, the only v1 host
encoding is
\begin{equation}
 \HostEncode(N,H)
 :=\mathsf{Hdr}_{72}\parallel N\parallel H,
 \qquad
 \ByteLen(\HostEncode(N,H))=72+\ByteLen(N)+\ByteLen(H).
 \label{eq:host-encoding-v1}
\end{equation}
Equation~\eqref{eq:host-encoding-v1} implies the four benchmark metrics
\begin{equation}
 \begin{aligned}
  \NargLen(\pi)&=\ByteLen(N),&
  \HintLen(\pi)&=\ByteLen(H),\\
  \PayloadLen(\pi)&=\ByteLen(N)+\ByteLen(H),&
  \WireLen(\pi)&=72+\PayloadLen(\pi).
 \end{aligned}
 \label{eq:wire-lengths-v1}
\end{equation}
The decoder \must{} verify Equation~\eqref{eq:wire-lengths-v1} with checked
arithmetic before allocating, enforce statement-derived upper bounds, and
require input length exactly $72+\mathtt{narg\_byte\_len}+
\mathtt{hint\_byte\_len}$.  The two declared record counts must match the
records actually consumed.  Header count and length fields are transport
metadata and are not transcript input.

\subsection{Typed event frames and registry}
\label{sec:wire-events}

Every NARG record, Hint record, public absorb, and squeeze event uses the same
prefix-free frame header:

\begin{lstlisting}[style=wire]
offset  size  field
0       2     event_tag                         u16-le
2       2     reserved = 0                     u16-le
4       4     chunk                            u32-le
8       4     level                            u32-le
12      4     round_or_index                   u32-le
16      8     payload_len                      u64-le
24      ...   payload[payload_len]
\end{lstlisting}

Let $\bot_{32}:=\mathtt{0xffffffff}$ denote an inapplicable scope coordinate.
For tag $\tau$, scope $(\ell,k,i)$, and payload $z$, the frame is exactly
\begin{equation}
 \Frame(\tau,\ell,k,i;z):=
 \LE{16}{\tau}\parallel\LE{16}{0}\parallel
 \LE{32}{\ell}\parallel\LE{32}{k}\parallel\LE{32}{i}\parallel
 \LE{64}{\ByteLen(z)}\parallel z.
 \label{eq:event-frame-v1}
\end{equation}
Thus Equation~\eqref{eq:event-frame-v1} has length $24+\ByteLen(z)$.  A stream
record with a reserved nonzero word, an inapplicable coordinate other than
$\bot_{32}$, a wrong channel, a wrong expected scope, or an unknown tag is
malformed.

The base-tag registry is below.  \textsf{D} events are public or derived,
have no proof record, and are absorbed at their ledger position;
\textsf{N} records occur only in $N$ and are absorbed;
\textsf{H} records occur only in $H$ and are checked but never absorbed;
\textsf{S} tags are squeeze requests.  For an \textsf{S} base tag $\tau$, the
squeeze-result tag is $\tau\mathbin{\mathrm{OR}}\mathtt{0x8000}$ and is derived,
never transmitted.  For an \textsf{S} row, the last column describes the
derived result; its request payload is always $\LE{32}{n}\parallel a$ as fixed
by Equation~\eqref{eq:transcript-v1}.

\begingroup\scriptsize
\renewcommand{\arraystretch}{1.04}
\setlength{\tabcolsep}{2pt}
\begin{longtable}{@{}p{0.52in}p{0.30in}p{1.62in}p{4.72in}@{}}
\toprule
\textbf{Tag}&\textbf{Class}&\textbf{Event}&\textbf{Scope; proof payload or squeeze result}\\
\midrule
\endfirsthead
\toprule
\textbf{Tag}&\textbf{Class}&\textbf{Event}&\textbf{Scope; proof payload or squeeze result}\\
\midrule
\endhead
\code{0001}&D&profile/domain&$(\bot,\bot,\bot)$; $\mathsf{ProfileV1}$ from
Equation~\eqref{eq:profile-statement-payloads}.\\
\code{0002}&D&statement&$(\bot,\bot,\bot)$; $\mathsf{StatementV1}$ from
Equation~\eqref{eq:profile-statement-payloads}.\\
\code{0003}&D&core seam&$(\bot,\bot,\bot)$;
$\mathsf{CoreSeamV1}$ from Equation~\eqref{eq:core-seam-payload}.\\
\code{1000--1fff}&profile&application PIOP&The selected application profile
owns a disjoint, byte-exact ledger within this range.\\
\code{2001}&N&chunk $\nu$&$(\ell,\bot,\bot)$; vector of $2^s$ elements of
$\K$.\\
\code{2081}&H&chunk $u$&$(\ell,\bot,\bot)$; vector of $2^s$ canonical
$u_{128}<2^{127}$.\\
\code{2101}&S&chunk $\zeta$&$(\ell,\bot,i)$ for $0\le i<s$; one $\K$
challenge.\\
\code{3001}&N&GKR phase-$c$ sum&$(\ell,k,\bot)$; one $\K$ element, absorbed
after the phase-link check and before the first phase-$c$ round.\\
\code{3002--30ff}&profile&GKR messages&The selected GKR profile owns exact
message tags, codecs, and checks in this range.\\
\code{3100--31ff}&profile&GKR squeezes&The selected GKR profile owns exact
squeeze tags and result types in this range.\\
\code{4001}&N&ring $s$ vector&$(\ell,\bot,\bot)$; vector of exactly 128
elements of $\K$.\\
\code{4101}&S&shared $r''$&$(\bot,\bot,i)$ for $0\le i<7$; one $\K$
challenge.\\
\code{4102}&S&batch $\eta$&$(\ell,\bot,\bot)$ in increasing active-chunk
order; one $\K$ challenge.\\
\code{5001}&D&opening target&$(\bot,\bot,\bot)$;
$\LE{32}{m_p}\parallel\operatorname{Enc}_{\K}(\beta)$.  $B$ is recomputed.\\
\code{5002}&N&initial root&$(\bot,0,\bot)$; one 32-byte SHA-256 digest.\\
\code{5003}&N&recursive root&$(\bot,k,\bot)$, $1\le k<J_m$; one digest.\\
\code{5004}&N&sumcheck message&scope by kind below; $\mathsf{kind}:u_8$
followed by $\operatorname{Enc}_{\K}(u_0)\parallel
\operatorname{Enc}_{\K}(u_2)$.\\
\code{5005}&N&OOD value&$(a,k,\bot)$; one $\K$ element for sample $a$.\\
\code{5006}&N&final $y_r$&$(\bot,J_m-1,\bot)$; vector of exactly
$2^{h_{J_m-1}}$ elements of $\K$.\\
\code{5007}&N&fold-PoW nonce&$(\bot,k,i)$; one $u_{64}$.\\
\code{5008}&N&query-PoW nonce&$(\bot,k,\bot)$; one $u_{64}$.\\
\code{5080}&H&initial opening&$(\bot,0,\bot)$;
Equation~\eqref{eq:merkle-opening-payload}.\\
\code{5081}&H&recursive opening&$(\bot,k,\bot)$, $1\le k<J_m-1$;
Equation~\eqref{eq:merkle-opening-payload}.\\
\code{5082}&H&final opening&$(\bot,J_m-1,\bot)$;
Equation~\eqref{eq:merkle-opening-payload}.\\
\code{5101}&S&fold challenge&$(\bot,k,i)$; one $\K$ challenge.\\
\code{5102}&S&OOD point&$(a,k,c)$; one $\K$ coordinate, with $c$ increasing.\\
\code{5103}&S&glue challenge&$(\bot,k,i)$ as fixed below; one $\K$ challenge.\\
\code{5104}&S&PoW seed&$(\bot,k,i)$; 16 bytes; request auxiliary payload $C$.\\
\code{5105}&S&query draw&$(\bot,k,a)$; one $\K$ challenge for attempt $a$.\\
\code{5106}&S&query alpha&$(\bot,k,c)$; one $\K$ challenge.\\
\bottomrule
\end{longtable}
\endgroup

The application and GKR ranges are black-box namespaces, not permission to
use implementation-defined bytes.  Their profile identifiers must select
published ledgers with the same columns as this table.  The fixed
\code{3001} event resolves the PCS-level phase-$c$ placement ambiguity in
Gap~G1; the rest of that gap remains subordinate-profile work and no profile
may reassign the tag.  The v1 outer pre-sumcheck
omits the redundant \code{presums[ell].claimed\_sums} field and derives
$[e_d^{(\ell)}+1]$, so tag \code{3081} is reserved and \mustnot{} occur.

\subsection{Transcript transitions and replay}
\label{sec:wire-transcript}

Model the transcript state as the byte string hashed by an unkeyed BLAKE3
hasher.  It starts empty.  The profile frame \code{0001} and statement frame
\code{0002} are absorbed, in that order, before any proof-dependent squeeze.
The application statement digest is
\begin{equation}
 d_{\rm app}:=\operatorname{BLAKE3}_{256}
 \bigl(\operatorname{hex}^{-1}
   (\code{46325a2d5043532f6170702d73746174656d656e742f763100})
 \parallel\mathsf{app\_profile\_id}\parallel
 \mathsf{app\_statement\_bytes}\bigr).
 \label{eq:app-statement-digest}
\end{equation}
The hex prefix in Equation~\eqref{eq:app-statement-digest} is the ASCII
literal \code{F2Z-PCS/app-statement/v1} followed by one NUL byte.  The two
initial frame payloads are exactly
\begin{equation}
 \begin{aligned}
  \mathsf{ProfileV1}:={}&
    \operatorname{hex}^{-1}(\code{46325a5043530000})
    \parallel\LE{16}{1}\parallel\ProfileID,\\
  \mathsf{StatementV1}:={}&
    \mathsf{app\_profile\_id}[32]\parallel
    \mathsf{gkr\_profile\_id}[32]\parallel
    \mathsf{code\_profile\_id}[32]\parallel d_{\rm app}[32]
    \parallel\LE{32}{t}\parallel\LE{32}{s}\parallel\LE{32}{W}\\
    &\parallel\operatorname{Enc}_{\K}(g)\parallel\rho[32].
 \end{aligned}
 \label{eq:profile-statement-payloads}
\end{equation}
Thus $\mathsf{ProfileV1}$ and $\mathsf{StatementV1}$ have 42 and 188 bytes,
respectively.  The application profile must make its final byte string in
Equation~\eqref{eq:app-statement-digest} canonical.

The derived core-seam payload is
\begin{equation}
 \mathsf{CoreSeamV1}:=
  \LE{32}{2^t}\parallel\prod_{b=0}^{2^t-1}\LE{128}{w_b}
  \parallel\LE{32}{2^s}\parallel
    \prod_{c=0}^{2^s-1}\LE{128}{\operatorname{can}_q(w'_c)}
  \parallel\LE{128}{\operatorname{can}_q(y)},
 \label{eq:core-seam-payload}
\end{equation}
with exactly $24+16(2^t+2^s)$ bytes.  It is absorbed under tag \code{0003}
after deterministic Bitify and before the first tag-\code{2001} record.

For a canonical frame $F$, absorption appends the exact frame bytes.  A
squeeze under base tag $\tau<\mathtt{0x8000}$, scope $(\ell,k,i)$, output
length $n$, and auxiliary request payload $a$ is the transition
\begin{equation}
 \begin{aligned}
  R&:=\Frame(\tau,\ell,k,i;\LE{32}{n}\parallel a),\\
  \tr'&:=\tr\parallel R,\\
  z&:=\operatorname{BLAKE3\text{-}XOF}(\tr')[0\mathbin{..}n],\\
  \tr''&:=\tr'\parallel
    \Frame(\tau\mathbin{\mathrm{OR}}\mathtt{0x8000},\ell,k,i;z).
 \end{aligned}
 \label{eq:transcript-v1}
\end{equation}
Equation~\eqref{eq:transcript-v1} uses a clone/finalize-XOF of the incremental
hasher after $R$; it does not reset or key the hasher.  A $\K$ challenge has
$n=16$ and decodes by Section~\ref{sec:wire-primitives}.  A vector challenge
is a sequence of scalar squeezes, not one long squeeze.  Unless its registry
row says otherwise, the scalars use consecutive $i$ values and the auxiliary
request payload $a$ is empty.  Tag \code{4102} instead uses
$(\ell,\bot,\bot)$ for each active $\ell$ in increasing order.  Tag
\code{5103} uses its one-byte glue kind, and tag \code{5104} uses the PoW
context $C$; these are the only nonempty core auxiliary payloads.

On a normal NARG read, the verifier validates the expected tag, scope, payload
length, and canonical payload, then immediately appends the complete frame
from Equation~\eqref{eq:event-frame-v1}.  A PoW nonce is the sole deferred
NARG read: the verifier parses its expected canonical frame, permits no
intervening event, checks Equation~\eqref{eq:pow-predicate-v1}, and appends the
frame only on success.  On a Hint read, it validates the same framing but
never appends those bytes.  It must complete the hint predicate before the
next squeeze.  In particular, it reads and absorbs \code{2001}, then reads
\code{2081} and verifies the exponent bounds and $g^u=\nu$ before the first
\code{2101} event.  NARG replay follows verifier-event order, never Rust
structure-field order.  Acceptance requires no deferred record or pending
hint and exact logical and byte exhaustion of both streams and every nested
slice.

\subsection{Grinding and query derivation}
\label{sec:wire-pow}

Let kind byte zero denote a fold grind and kind byte one a query grind.  For
kind $K_{\rm pow}$, level $k$, round $i$, and derived difficulty $d$, define
the exact PoW context
\begin{equation}
 C:=\operatorname{ASCII}(\code{F2Z-PCS/PoW/v1})\parallel
 \LE{8}{K_{\rm pow}}\parallel\LE{32}{k}\parallel
 \LE{32}{i}\parallel\LE{16}{d}.
 \label{eq:pow-context-v1}
\end{equation}
In Equation~\eqref{eq:pow-context-v1}, $i=\bot_{32}$ for a query grind.  The seed is the 16-byte output of
Equation~\eqref{eq:transcript-v1} under tag \code{5104}, the same scope, and
auxiliary request payload $C$.  After obtaining the seed, the verifier reads
the expected \code{5007} or \code{5008} NARG frame without absorbing it and
requires, for the scheduled positive difficulty $d$,
\begin{equation}
 \LZ\!\left(\operatorname{BLAKE3}_{256}
   (C\parallel\mathsf{seed}_{16}\parallel\LE{64}{n})\right)\ge d,
 \qquad d>0.
 \label{eq:pow-predicate-v1}
\end{equation}
$\LZ$ in Equation~\eqref{eq:pow-predicate-v1} counts zero bits from the
most-significant bit of digest byte zero, then proceeds toward the
least-significant bit and subsequent bytes.  Any satisfying nonce is accepted;
minimality is not checked.  Only after the predicate succeeds does
the verifier absorb the complete nonce frame, before any dependent squeeze.
Difficulty zero has no predicate invocation, seed squeeze, or nonce record in
this profile.

The exact scheduled difficulties are
\begin{equation}
 d^{\rm query}_k:=g_k=16,
 \qquad
 d^{\rm fold}_{k,i}:=\max(f_k-i,0).
 \label{eq:pow-difficulties-v1}
\end{equation}
Equation~\eqref{eq:pow-difficulties-v1} is evaluated with ordinary
nonnegative integers.  After the preceding sumcheck message and immediately
before fold challenge $(k,i)$, perform both the tag-\code{5104} seed squeeze
and tag-\code{5007} nonce record if $d^{\rm fold}_{k,i}>0$; if the difficulty
is zero, perform neither event.  For $0\le k<J_m-1$, query event $k$ occurs after the
tag-\code{5003} commitment root for level $k+1$, all $o_{k+1}$ OOD bindings
for that root, and their glue challenges, and immediately before the tag-\code{5105}
draws that open level $k$.  Query event $J_m-1$ instead occurs after the
tag-\code{5006} final $y_r$, with no new root or OOD binding, and immediately
before the final-level draws.  There is exactly one query nonce at every level.
Consequently the exact number of nonce records is
\begin{equation}
 N_{\rm pow}=J_m+\sum_{k=0}^{J_m-1}\min(k_k,f_k).
 \label{eq:pow-count-v1}
\end{equation}
Missing, extra, reordered, or trailing nonces violate
Equation~\eqref{eq:pow-count-v1} and reject.

For level $k$, each \code{5105} event yields $r\in\K$ and proposes
\begin{equation}
 q:=r_{\lo}\bmod 2^{h_k+r_k}.
 \label{eq:query-position-v1}
\end{equation}
The verifier repeats Equation~\eqref{eq:query-position-v1} until $Q_k$
distinct positions exist, discarding duplicates, and then sorts the positions
in increasing order.  It then derives the tag-\code{5106} alphas before
reading and checking the opening.

\subsection{Recursive-opening replay ledger}
\label{sec:wire-recursive-ledger}

This subsection fixes the complete replay order for the recursive opening.
The four canonical tag-\code{5004} message kinds and scopes are
\begin{equation}
 \begin{array}{ccll}
  \mathtt{0x00}&\mathsf{Start}&(\bot,0,\bot)&\text{exactly once},\\
  \mathtt{0x01}&\mathsf{FoldNext}&(\bot,k,j)&0\le j<k_k,\\
  \mathtt{0x02}&\mathsf{OodIntro}&(a,k,\bot)&1\le k<J_m, 0\le a<o_k,\\
  \mathtt{0x03}&\mathsf{QueryIntro}&(\bot,k,\bot)&0\le k<J_m-1.
 \end{array}
 \label{eq:sumcheck-message-kinds}
\end{equation}
Each payload is the one-byte kind followed by the two canonical $\K$
elements in the registry.  Any other kind/scope pair rejects.

The outer $\mathsf{VerifyF2Z}$ caller first requires that $B$ is the
deterministic coefficient function derived in Group~9, with domain size
$2^{m_p}$; $B$ is never densely serialized.  The caller alone absorbs tag
\code{5001} with payload
$\LE{32}{m_p}\parallel\operatorname{Enc}_{\K}(\beta)$ and then invokes
$\mathsf{RecVerifyStream}$.  That helper starts after tag \code{5001}: its
first event reads tag \code{5002}, requires the root to equal $\rho$, absorbs
it, and then reads and absorbs the $\mathsf{Start}$ tag-\code{5004} message.

For each $k=0,\ldots,J_m-1$ in increasing order, it then performs the
following ledger without interleaving another event.
\begin{enumerate}
 \item For $j=0,\ldots,k_k-1$, perform the optional fold grind exactly as in
 Section~\ref{sec:wire-pow}, squeeze tag \code{5101} at
 $(\bot,k,j)$, then read, validate, and absorb the $\mathsf{FoldNext}$
 tag-\code{5004} message at the same scope.

 \item If $k<J_m-1$, read and absorb the tag-\code{5003} root at
 $(\bot,k+1,\bot)$.  For each $a=0,\ldots,o_{k+1}-1$, squeeze the
 $n_{k+1}$ tag-\code{5102} coordinates at scopes $(a,k+1,c)$ for
 $c=0,\ldots,n_{k+1}-1$; read and absorb tag \code{5005} at
 $(a,k+1,\bot)$; read and absorb the $\mathsf{OodIntro}$ tag-\code{5004}
 message at that scope; and squeeze tag \code{5103} at
 $(\bot,k+1,a)$ with the one-byte auxiliary kind \code{00}.

 If instead $k=J_m-1$, read and absorb tag \code{5006} at
 $(\bot,k,\bot)$, requiring exactly $2^{h_k}$ canonical $\K$ elements.

 \item Perform the query grind at $(\bot,k,\bot)$ with difficulty 16 and
 absorb its validated tag-\code{5008} nonce.  Starting with attempt $a=0$,
 squeeze tag \code{5105} at $(\bot,k,a)$ and increment $a$ after every
 squeeze, including a duplicate, until $Q_k$ distinct positions have been
 obtained.  Reject rather than wrap if the next attempt would exceed
 $2^{32}-1$.  Sort the positions, then squeeze exactly
 $\lceil\log_2Q_k\rceil$ tag-\code{5106} alphas at $(\bot,k,c)$ for
 consecutive $c$.

 Read exactly one opening Hint using tag \code{5080} for $k=0$, tag
 \code{5081} for $0<k<J_m-1$, or tag \code{5082} for $k=J_m-1$.  Check
 Equations~\eqref{eq:merkle-opening-payload}--\eqref{eq:merkle-opening-length}
 and the root before any later event.  For $k<J_m-1$, next read and absorb
 the $\mathsf{QueryIntro}$ tag-\code{5004} message.  At every level,
 including the final level where no such message exists, squeeze the query
 glue tag \code{5103} at $(\bot,k,o_k)$ with one-byte auxiliary kind
 \code{01}.  At the final level this challenge is used directly in the
 residual check.
\end{enumerate}

The fixed record/event cardinalities are therefore
\begin{equation}
 \begin{aligned}
  \#5002&=1,& \#5003&=J_m-1,& \#5005&=\sum_{k=1}^{J_m-1}o_k,\\
  \#5004&=1+\sum_{k=0}^{J_m-1}k_k
     +\sum_{k=1}^{J_m-1}o_k+(J_m-1),&
  \#5006&=1,& \#508{*}&=J_m,\\
  \#5101&=\sum_{k=0}^{J_m-1}k_k,&
  \#5102&=\sum_{k=1}^{J_m-1}o_kn_k,&
  \#5103&=J_m+\sum_{k=1}^{J_m-1}o_k,\\
  \#5106&=\sum_{k=0}^{J_m-1}\lceil\log_2Q_k\rceil.
 \end{aligned}
 \label{eq:recursive-event-counts}
\end{equation}
Here $\#508{*}$ is one opening Hint per level.  Tag-\code{5105} count is the
transcript-derived number of attempts, all of which appear in the audit log.
Nonce counts remain Equation~\eqref{eq:pow-count-v1}.  A different order,
scope, kind, cardinality, or unread field rejects.

\subsection{Parser, implementation, and conformance requirements}
\label{sec:wire-conformance}

A v1 implementation \must{}:
\begin{enumerate}
 \item validate magic, version, flags, profile and subordinate profile
 identifiers, all profile-, statement-, and transcript-derived bounds, and all
 reserved fields before allocating or doing proof work;
 \item use checked arithmetic and checked conversion for every length and
 index, reject noncanonical field values, and return a decoding error rather
 than panic on adversarial bytes;
 \item encode the NARG stream in the exact absorb order and the Hint stream in
 its exact check order, with independent cursors;
 \item replace opaque recursive-proof serialization with the field records and
 exact order in Sections~\ref{sec:wire-events}
 and~\ref{sec:wire-recursive-ledger}; and
 \item reject unless every outer, subordinate, recursive, row, and Merkle-path
 slice is exactly exhausted.
\end{enumerate}

Conformance fixtures \must{} include canonical primitive and frame hex, one
complete $N/H/\HostEncode$ proof, the ordered transcript audit log, every
squeezed challenge, PoW seeds/nonces and query positions, and all four lengths
from Equation~\eqref{eq:wire-lengths-v1}.  Negative fixtures \must{} cover an
unknown version/profile/tag, a reserved bit, wrong endian or channel, record
reordering, truncation and suffix bytes, noncanonical scalars, count overflow,
a failed or late $u/\nu$ hint, altered PoW scope/difficulty/nonce, and missing
or extra recursive fields.  Native, recursive, and independent reference
verifiers must produce byte-identical transcript audit logs.

The repository file \code{wire-v1-vectors.txt} is the normative initial
known-answer set for the profile identifier, primitive container/frame codec,
Merkle row/tree/multiproof codec, and PoW predicate at difficulties 1, 7, 8,
and 16.  A conforming implementation must reproduce it exactly and add a
complete proof/transcript
fixture before claiming end-to-end v1 interoperability.

\endgroup
